% answers.tex —— 2013 年·全国大纲卷·理科 · 参考答案与详细解析片段
% 由 tools/generate_projects.py 依来源总账自 2013/2013_数学真题_大纲理_试题.md + 2013/2013_数学真题_大纲理_解析正文.md + 2013/2013_数学真题_大纲理_答案要点.md 生成，勿手改（改源 MD 后重跑）


\begin{answer}
\gksolhead{1}{}

\textbf{题目}：

设集合 $A=\{1,2,3\}$，$B=\{4,5\}$，$M=\{x\mid x=a+b,\ a\in A,\ b\in B\}$，则 $M$ 中元素的个数为\mbox{（　　）}

（A）$3$；（B）$4$；（C）$5$；（D）$6$。

\Needspace{4\baselineskip}
\textbf{解析}
\Needspace{5\baselineskip}
$A=\{1,2,3\}$，$B=\{4,5\}$，逐对相加：
\nopagebreak[4]
\[1+4=5,\quad 1+5=6,\quad 2+4=6,\quad 2+5=7,\quad 3+4=7,\quad 3+5=8.\]
\Needspace{5\baselineskip}
故
\nopagebreak[4]
\[M=\{5,6,7,8\},\]
共有 $4$ 个元素。选 B。

\gkanswer{答案：B。}

\end{answer}
\begin{answer}
\gksolhead{2}{}

\textbf{题目}：

$(1+\sqrt3\mathrm i)^3=$\mbox{（　　）}

（A）$-8$；（B）$8$；（C）$-8\mathrm i$；（D）$8\mathrm i$。

\Needspace{4\baselineskip}
\textbf{解析}
\Needspace{5\baselineskip}
先平方：
\nopagebreak[4]
\[(1+\sqrt3\mathrm i)^2=1+2\sqrt3\mathrm i+3\mathrm i^2
=-2+2\sqrt3\mathrm i.\]
\Needspace{5\baselineskip}
再乘 $1+\sqrt3\mathrm i$：
\nopagebreak[4]
\[\begin{aligned}
(1+\sqrt3\mathrm i)^3
&=(-2+2\sqrt3\mathrm i)(1+\sqrt3\mathrm i)\\
&=-2-2\sqrt3\mathrm i+2\sqrt3\mathrm i+2\sqrt3\mathrm i\cdot\sqrt3\mathrm i\\
&=-2-6=-8.
\end{aligned}\]
选 A。

\gkanswer{答案：A。}

\end{answer}
\begin{answer}
\gksolhead{3}{}

\textbf{题目}：

已知向量 $\boldsymbol m=(\lambda+1,1)$，$\boldsymbol n=(\lambda+2,2)$，若 $(\boldsymbol m+\boldsymbol n)\perp(\boldsymbol m-\boldsymbol n)$，则 $\lambda=$\mbox{（　　）}

（A）$-4$；（B）$-3$；（C）$-2$；（D）$-1$。

\Needspace{4\baselineskip}
\textbf{解析}
\Needspace{5\baselineskip}
由 $(\boldsymbol m+\boldsymbol n)\perp(\boldsymbol m-\boldsymbol n)$ 得
\nopagebreak[4]
\[(\boldsymbol m+\boldsymbol n)\cdot(\boldsymbol m-\boldsymbol n)
=|\boldsymbol m|^2-|\boldsymbol n|^2=0.\]
\Needspace{5\baselineskip}
而
\nopagebreak[4]
\[|\boldsymbol m|^2=(\lambda+1)^2+1=\lambda^2+2\lambda+2,\]
\nopagebreak[4]
\[|\boldsymbol n|^2=(\lambda+2)^2+4=\lambda^2+4\lambda+8.\]
\Needspace{5\baselineskip}
于是
\nopagebreak[4]
\[\lambda^2+2\lambda+2=\lambda^2+4\lambda+8,\]
即 $2\lambda=-6$，故 $\lambda=-3$。选 B。

\gkanswer{答案：B。}

\end{answer}
\begin{answer}
\gksolhead{4}{}

\textbf{题目}：

已知函数 $f(x)$ 的定义域为 $(-1,0)$，则函数 $f(2x+1)$ 的定义域为\mbox{（　　）}

（A）$(-1,1)$；（B）$\left(-1,-\dfrac12\right)$；（C）$(-1,0)$；（D）$\left(\dfrac12,1\right)$。

\Needspace{4\baselineskip}
\textbf{解析}
\Needspace{5\baselineskip}
$f$ 的定义域为 $(-1,0)$，故 $f(2x+1)$ 有意义当且仅当
\nopagebreak[4]
\[-1<2x+1<0.\]
\Needspace{5\baselineskip}
各边减 $1$：
\nopagebreak[4]
\[-2<2x<-1,\]
\Needspace{5\baselineskip}
再除以 $2$：
\nopagebreak[4]
\[-1<x<-\frac12.\]
所以定义域为 $\left(-1,-\dfrac12\right)$。选 B。

\gkanswer{答案：B。}

\end{answer}
\begin{answer}
\gksolhead{5}{}

\textbf{题目}：

函数 $f(x)=\log_2\left(1+\dfrac1x\right)$（$x>0$）的反函数 $f^{-1}(x)=$\mbox{（　　）}

（A）$\dfrac1{2^x-1}$（$x>0$）；（B）$\dfrac1{2^x-1}$（$x\neq0$）；（C）$2^x-1$（$x\in\mathbf R$）；（D）$2^x-1$（$x>0$）。

\Needspace{4\baselineskip}
\textbf{解析}
\Needspace{5\baselineskip}
设 $y=f(x)=\log_2\left(1+\dfrac1x\right)$，则
\nopagebreak[4]
\[2^y=1+\frac1x
\Longrightarrow \frac1x=2^y-1
\Longrightarrow x=\frac1{2^y-1}.\]
\Needspace{5\baselineskip}
当 $x>0$ 时 $\dfrac1x>0$，故 $1+\dfrac1x>1$，从而 $y>0$。交换 $x$，$y$，得
\nopagebreak[4]
\[f^{-1}(x)=\frac1{2^x-1}\quad(x>0).\]
选 A。

\gkanswer{答案：A。}

\end{answer}
\begin{answer}
\gksolhead{6}{}

\textbf{题目}：

已知数列 $\{a_n\}$ 满足 $3a_{n+1}+a_n=0$，$a_2=-\dfrac43$，则 $\{a_n\}$ 的前 $10$ 项和等于\mbox{（　　）}

（A）$-6(1-3^{-10})$；（B）$\dfrac19(1-3^{10})$；（C）$3(1-3^{-10})$；（D）$3(1+3^{-10})$。

\Needspace{4\baselineskip}
\textbf{解析}
\Needspace{5\baselineskip}
由 $3a_{n+1}+a_n=0$ 得
\nopagebreak[4]
\[a_{n+1}=-\frac13a_n,\]
\Needspace{5\baselineskip}
故 $\{a_n\}$ 是公比 $q=-\dfrac13$ 的等比数列。由
\nopagebreak[4]
\[a_2=a_1q=-\frac{a_1}3=-\frac43\]
\Needspace{5\baselineskip}
得 $a_1=4$。前 $10$ 项和为
\nopagebreak[4]
\[\begin{aligned}
S_{10}
&=\frac{4\left[1-\left(-\frac13\right)^{10}\right]}{1-\left(-\frac13\right)}
=\frac{4\left(1-3^{-10}\right)}{\frac43}\\
&=3\left(1-3^{-10}\right).
\end{aligned}\]
选 C。

\gkanswer{答案：C。}

\end{answer}
\begin{answer}
\gksolhead{7}{}

\textbf{题目}：

$(1+x)^8(1+y)^4$ 的展开式中 $x^2y^2$ 的系数是\mbox{（　　）}

（A）$56$；（B）$84$；（C）$112$；（D）$168$。

\Needspace{4\baselineskip}
\textbf{解析}
\Needspace{5\baselineskip}
$(1+x)^8$ 中 $x^2$ 的系数为 $\mathrm C_8^2=28$，$(1+y)^4$ 中 $y^2$ 的系数为 $\mathrm C_4^2=6$。两展开式相乘，$x^2y^2$ 的系数为
\nopagebreak[4]
\[\mathrm C_8^2\mathrm C_4^2=28\times6=168.\]
选 D。

\gkanswer{答案：D。}

\end{answer}
\begin{answer}
\gksolhead{8}{}

\textbf{题目}：

椭圆 $C:\dfrac{x^2}4+\dfrac{y^2}3=1$ 的左、右顶点分别为 $A_1$、$A_2$，点 $P$ 在 $C$ 上且直线 $PA_2$ 斜率的取值范围是 $[-2,-1]$，那么直线 $PA_1$ 斜率的取值范围是\mbox{（　　）}

（A）$\left[\dfrac12,\dfrac34\right]$；（B）$\left[\dfrac38,\dfrac34\right]$；（C）$\left[\dfrac12,1\right]$；（D）$\left[\dfrac34,1\right]$。

\Needspace{4\baselineskip}
\textbf{解析}
\Needspace{5\baselineskip}
椭圆中 $a=2$，$b=\sqrt3$，左、右顶点为 $A_1(-2,0)$，$A_2(2,0)$。设 $P(x_0,y_0)$ 在椭圆上且不是顶点，则 $y_0\neq0$，$x_0\neq\pm2$。两条直线的斜率分别为
\nopagebreak[4]
\[k_{PA_1}=\frac{y_0}{x_0+2},\qquad
k_{PA_2}=\frac{y_0}{x_0-2}.\]
\Needspace{5\baselineskip}
相乘：
\nopagebreak[4]
\[k_{PA_1}k_{PA_2}
=\frac{y_0^2}{x_0^2-4}.\]
\Needspace{5\baselineskip}
由椭圆方程
\nopagebreak[4]
\[y_0^2=3\left(1-\frac{x_0^2}4\right)
=\frac34(4-x_0^2),\]
\Needspace{5\baselineskip}
故
\nopagebreak[4]
\[k_{PA_1}k_{PA_2}
=\frac{\frac34(4-x_0^2)}{x_0^2-4}
=-\frac34.\]
\Needspace{5\baselineskip}
于是
\nopagebreak[4]
\[k_{PA_1}=-\frac{3}{4k_{PA_2}}.\]
\Needspace{5\baselineskip}
函数 $k_{PA_1}=-\dfrac{3}{4k_{PA_2}}$ 在 $k_{PA_2}<0$ 时单调递增，因为其导数为 $\dfrac{3}{4k_{PA_2}^2}>0$。当 $k_{PA_2}=-2$ 时 $k_{PA_1}=\dfrac38$；当 $k_{PA_2}=-1$ 时 $k_{PA_1}=\dfrac34$。所以
\nopagebreak[4]
\[k_{PA_1}\in\left[\frac38,\frac34\right].\]
选 B。

\gkanswer{答案：B。}

\end{answer}
\begin{answer}
\gksolhead{9}{}

\textbf{题目}：

若函数 $f(x)=x^2+ax+\dfrac1x$ 在 $\left(\dfrac12,+\infty\right)$ 是增函数，则 $a$ 的取值范围是\mbox{（　　）}

（A）$[-1,0]$；（B）$[-1,+\infty)$；（C）$[0,3]$；（D）$[3,+\infty)$。

\Needspace{4\baselineskip}
\textbf{解析}
\Needspace{5\baselineskip}
求导：
\nopagebreak[4]
\[f'(x)=2x+a-\frac1{x^2}.\]
\Needspace{5\baselineskip}
$f(x)$ 在 $\left(\dfrac12,+\infty\right)$ 上为增函数，需且只需
\nopagebreak[4]
\[f'(x)\geqslant0\quad\text{\text{对一切} }x>\frac12\text{ \text{成立}}.\]
\Needspace{5\baselineskip}
即
\nopagebreak[4]
\[a\geqslant\frac1{x^2}-2x\quad\left(x>\frac12\right).\]
\Needspace{5\baselineskip}
记
\nopagebreak[4]
\[g(x)=\frac1{x^2}-2x,\]
\Needspace{5\baselineskip}
则
\nopagebreak[4]
\[g'(x)=-\frac2{x^3}-2<0\quad\left(x>\frac12\right),\]
\Needspace{5\baselineskip}
故 $g(x)$ 在 $\left(\dfrac12,+\infty\right)$ 上单调递减，从而
\nopagebreak[4]
\[g(x)<g\left(\frac12\right)=4-1=3.\]
\Needspace{5\baselineskip}
因此 $a\geqslant3$ 时不等式对一切 $x>\dfrac12$ 成立；而 $a<3$ 时，取 $x>\dfrac12$ 充分接近 $\dfrac12$，仍有 $g(x)>a$，不满足。故
\nopagebreak[4]
\[a\in[3,+\infty).\]
选 D。

\gkanswer{答案：D。}

\end{answer}
\begin{answer}
\gksolhead{10}{}

\textbf{题目}：

已知正四棱柱 $ABCD\text{-}A_1B_1C_1D_1$ 中，$AA_1=2AB$，则 $CD$ 与平面 $BDC_1$ 所成角的正弦值等于\mbox{（　　）}

（A）$\dfrac23$；（B）$\dfrac{\sqrt3}3$；（C）$\dfrac{\sqrt2}3$；（D）$\dfrac13$。

\Needspace{4\baselineskip}
\textbf{解析}
\Needspace{5\baselineskip}
设 $AB=1$，则 $AA_1=2$。以 $A$ 为原点，$AB$，$AD$，$AA_1$ 所在直线分别为 $x$，$y$，$z$ 轴建立坐标系：
\nopagebreak[4]
\[A(0,0,0),\quad B(1,0,0),\quad D(0,1,0),\]
\nopagebreak[4]
\[C(1,1,0),\quad C_1(1,1,2).\]
\Needspace{5\baselineskip}
平面 $BDC_1$ 内两向量
\nopagebreak[4]
\[\overrightarrow{DB}=(1,-1,0),\qquad
\overrightarrow{DC_1}=(1,0,2).\]
\Needspace{5\baselineskip}
其法向量可取
\nopagebreak[4]
\[\boldsymbol n=\overrightarrow{DB}\times\overrightarrow{DC_1}
=(-2,-2,1),\qquad|\boldsymbol n|=3.\]
\Needspace{5\baselineskip}
又 $CD$ 的方向向量
\nopagebreak[4]
\[\overrightarrow{CD}=(-1,0,0).\]
\Needspace{5\baselineskip}
设 $CD$ 与平面 $BDC_1$ 所成角为 $\theta$，则
\nopagebreak[4]
\[\sin\theta
=\frac{\left|\overrightarrow{CD}\cdot\boldsymbol n\right|}
{|\overrightarrow{CD}|\,|\boldsymbol n|}
=\frac{2}{1\times3}=\frac23.\]
选 A。

\gkanswer{答案：A。}

\end{answer}
\begin{answer}
\gksolhead{11}{}

\textbf{题目}：

已知抛物线 $C:y^2=8x$ 与点 $M(-2,2)$，过 $C$ 的焦点且斜率为 $k$ 的直线与 $C$ 交于 $A$、$B$ 两点，若 $\overrightarrow{MA}\cdot\overrightarrow{MB}=0$，则 $k=$\mbox{（　　）}

（A）$\dfrac12$；（B）$\dfrac{\sqrt2}2$；（C）$\sqrt2$；（D）$2$。

\Needspace{4\baselineskip}
\textbf{解析}
\Needspace{5\baselineskip}
抛物线 $y^2=8x$ 的焦点为 $F(2,0)$。过 $F$ 且斜率为 $k$ 的直线为
\nopagebreak[4]
\[y=k(x-2)\quad(k\neq0),\]
\Needspace{5\baselineskip}
即 $x=2+\dfrac yk$。代入 $y^2=8x$：
\nopagebreak[4]
\[y^2=8\left(2+\frac yk\right),\]
\Needspace{5\baselineskip}
整理得
\nopagebreak[4]
\[y^2-\frac8k y-16=0.\]
\Needspace{5\baselineskip}
设 $A(x_1,y_1)$，$B(x_2,y_2)$，由韦达定理
\nopagebreak[4]
\[y_1+y_2=\frac8k,\qquad y_1y_2=-16.\]
\Needspace{5\baselineskip}
且 $x_i=2+\dfrac{y_i}{k}$。由 $\overrightarrow{MA}\cdot\overrightarrow{MB}=0$ 得
\nopagebreak[4]
\[(x_1+2)(x_2+2)+(y_1-2)(y_2-2)=0.\]
\Needspace{5\baselineskip}
先算第一部分：
\nopagebreak[4]
\[\begin{aligned}
(x_1+2)(x_2+2)
&=\left(4+\frac{y_1}{k}\right)\left(4+\frac{y_2}{k}\right)\\
&=16+\frac{4(y_1+y_2)}{k}+\frac{y_1y_2}{k^2}\\
&=16+\frac{32}{k^2}-\frac{16}{k^2}
=16+\frac{16}{k^2}.
\end{aligned}\]
\Needspace{5\baselineskip}
再算第二部分：
\nopagebreak[4]
\[\begin{aligned}
(y_1-2)(y_2-2)
&=y_1y_2-2(y_1+y_2)+4\\
&=-16-\frac{16}{k}+4
=-12-\frac{16}{k}.
\end{aligned}\]
\Needspace{5\baselineskip}
代入点积条件：
\nopagebreak[4]
\[16+\frac{16}{k^2}-12-\frac{16}{k}=0,\]
\Needspace{5\baselineskip}
即
\nopagebreak[4]
\[4+\frac{16}{k^2}-\frac{16}{k}=0.\]
\Needspace{5\baselineskip}
两边乘以 $k^2$：
\nopagebreak[4]
\[4k^2-16k+16=0
\Longrightarrow (k-2)^2=0,\]
故 $k=2$。选 D。

\gkanswer{答案：D。}

\end{answer}
\begin{answer}
\gksolhead{12}{}

\textbf{题目}：

已知函数 $f(x)=\cos x\sin2x$，下列结论中错误的是\mbox{（　　）}

（A）$y=f(x)$ 的图象关于点 $(\pi,0)$ 中心对称；（B）$y=f(x)$ 的图象关于 $x=\dfrac\pi2$ 对称；（C）$f(x)$ 的最大值为 $\dfrac{\sqrt3}2$；（D）$f(x)$ 既是奇函数，又是周期函数。

\Needspace{4\baselineskip}
\textbf{解析}
\Needspace{5\baselineskip}
利用 $\sin2x=2\sin x\cos x$，有
\nopagebreak[4]
\[f(x)=2\sin x\cos^2x
=2\sin x(1-\sin^2x)
=2(\sin x-\sin^3x).\]
逐条检验。

\Needspace{5\baselineskip}
对于（A），因为
\nopagebreak[4]
\[f(2\pi-x)=\cos(2\pi-x)\sin(4\pi-2x)
=\cos x(-\sin2x)=-f(x),\]
所以图象关于点 $(\pi,0)$ 中心对称，正确。

\Needspace{5\baselineskip}
对于（B），因为
\nopagebreak[4]
\[f(\pi-x)=\cos(\pi-x)\sin(2\pi-2x)
=(-\cos x)(-\sin2x)=f(x),\]
所以图象关于 $x=\dfrac\pi2$ 对称，正确。

\Needspace{5\baselineskip}
对于（C），令 $t=\sin x\in[-1,1]$，则
\nopagebreak[4]
\[f(x)=2(t-t^3)=:h(t).\]
\Needspace{5\baselineskip}
求导
\nopagebreak[4]
\[h'(t)=2(1-3t^2),\]
\Needspace{5\baselineskip}
故 $h$ 在 $t=\pm\dfrac1{\sqrt3}$ 处取得极值，且
\nopagebreak[4]
\[h\left(\frac1{\sqrt3}\right)
=2\left(\frac1{\sqrt3}-\frac1{3\sqrt3}\right)
=\frac4{3\sqrt3}
=\frac{4\sqrt3}{9}.\]
\Needspace{5\baselineskip}
端点值 $h(\pm1)=0$，所以最大值为 $\dfrac{4\sqrt3}{9}$。而
\nopagebreak[4]
\[\frac{4\sqrt3}{9}\neq\frac{\sqrt3}{2},\]
故（C）错误。

对于（D），$f(-x)=\cos x\sin(-2x)=-f(x)$，且 $f(x+2\pi)=f(x)$，所以既是奇函数又是周期函数，正确。

因此错误的是（C）。选 C。

\gkanswer{答案：C。}

\end{answer}
\begin{answer}
\gksolhead{13}{}

\textbf{题目}：

已知 $\alpha$ 是第三象限角，$\sin\alpha=-\dfrac13$，则 $\cot\alpha=$\gkblank{}。

\Needspace{4\baselineskip}
\textbf{解析}
\Needspace{5\baselineskip}
$\alpha$ 是第三象限角，故 $\cos\alpha<0$。由
\nopagebreak[4]
\[\sin^2\alpha+\cos^2\alpha=1\]
\Needspace{5\baselineskip}
得
\nopagebreak[4]
\[\cos^2\alpha=1-\left(-\frac13\right)^2
=1-\frac19=\frac89,\]
\Needspace{5\baselineskip}
所以
\nopagebreak[4]
\[\cos\alpha=-\frac{2\sqrt2}{3}.\]
\Needspace{5\baselineskip}
于是
\nopagebreak[4]
\[\cot\alpha=\frac{\cos\alpha}{\sin\alpha}
=\frac{-\frac{2\sqrt2}{3}}{-\frac13}
=2\sqrt2.\]

\gkanswer{答案：$2\sqrt2$。}

\end{answer}
\begin{answer}
\gksolhead{14}{}

\textbf{题目}：

$6$ 个人排成一排，其中甲、乙两人不相邻的不同排法共有\gkblank{}种。（用数字作答）

\Needspace{4\baselineskip}
\textbf{解析}
\Needspace{5\baselineskip}
$6$ 人任意排列共有 $6!$ 种。甲、乙相邻时，把两人看成一个整体，与其余 $4$ 人共 $5$ 个对象排列，再考虑甲、乙内部的顺序，共有
\nopagebreak[4]
\[2\cdot5!\]
\Needspace{5\baselineskip}
种。故甲、乙不相邻的排法数为
\nopagebreak[4]
\[6!-2\cdot5!=720-240=480.\]

\gkanswer{答案：$480$。}

\end{answer}
\begin{answer}
\gksolhead{15}{}

\textbf{题目}：

记不等式组 $\begin{cases}x\geqslant0,\\ x+3y\geqslant4,\\ 3x+y\leqslant4\end{cases}$ 所表示的平面区域为 $D$，若直线 $y=a(x+1)$ 与 $D$ 有公共点，则 $a$ 的取值范围是\gkblank{}。

\Needspace{4\baselineskip}
\textbf{解析}
\Needspace{5\baselineskip}
三条边界直线为
\nopagebreak[4]
\[x=0,\qquad x+3y=4,\qquad 3x+y=4.\]
\Needspace{5\baselineskip}
两两求交点：
\nopagebreak[4]
\[x=0\text{ \text{与} }x+3y=4\Longrightarrow \left(0,\frac43\right),\]
\nopagebreak[4]
\[x=0\text{ \text{与} }3x+y=4\Longrightarrow (0,4),\]
\nopagebreak[4]
\[x+3y=4\text{ \text{与} }3x+y=4\Longrightarrow (1,1).\]
\Needspace{5\baselineskip}
由 $x\geqslant0$ 及
\nopagebreak[4]
\[\frac{4-x}{3}\leqslant y\leqslant4-3x\]
\Needspace{5\baselineskip}
知 $0\leqslant x\leqslant1$，故可行域 $D$ 是以
\nopagebreak[4]
\[\left(0,\frac43\right),\qquad(1,1),\qquad(0,4)\]
为顶点的三角形。

\Needspace{5\baselineskip}
直线 $y=a(x+1)$ 过定点 $(-1,0)$，斜率为 $a$。该点与三角形三个顶点连线的斜率分别为
\nopagebreak[4]
\[\frac{\frac43-0}{0-(-1)}=\frac43,\qquad
\frac{1-0}{1-(-1)}=\frac12,\qquad
\frac{4-0}{0-(-1)}=4.\]
\Needspace{5\baselineskip}
其中最小值为 $\dfrac12$，最大值为 $4$。由于 $D$ 是凸集且全部位于 $x\geqslant0>-1$，过 $(-1,0)$ 的直线与 $D$ 有公共点当且仅当其斜率介于这两条边界斜率之间，故
\nopagebreak[4]
\[a\in\left[\frac12,4\right].\]

\gkanswer{答案：$\left[\dfrac12,4\right]$。}

\end{answer}
\begin{answer}
\gksolhead{16}{}

\textbf{题目}：

已知圆 $O$ 和圆 $K$ 是球 $O$ 的大圆和小圆，其公共弦长等于球 $O$ 的半径，$OK=\dfrac32$，且圆 $O$ 与圆 $K$ 所在的平面所成的一个二面角为 $60^\circ$，则球 $O$ 的表面积等于\gkblank{}。

\Needspace{4\baselineskip}
\textbf{解析}
\Needspace{5\baselineskip}
设球 $O$ 的半径为 $R$，圆 $O$ 为大圆，圆 $K$ 为小圆，两圆所在平面的公共弦为 $AB$，其中点为 $M$。则
\nopagebreak[4]
\[AB=R,\qquad OM\perp AB,\qquad KM\perp AB.\]
\Needspace{5\baselineskip}
在大圆 $O$ 中，弦心距
\nopagebreak[4]
\[OM=\sqrt{R^2-\left(\frac R2\right)^2}
=\frac{\sqrt3R}{2}.\]
又 $OK$ 垂直于小圆 $K$ 所在平面，而 $AB$ 在该平面内，故 $OK\perp AB$；于是 $OK\perp KM$，三角形 $OKM$ 在 $K$ 处为直角。

\Needspace{5\baselineskip}
两圆所在平面的二面角等于 $OM$ 与 $KM$ 的夹角。由 $OK\perp KM$ 知 $\angle OMK$ 为锐角，故
\nopagebreak[4]
\[\angle OMK=60^\circ.\]
\Needspace{5\baselineskip}
在直角三角形 $OMK$ 中，
\nopagebreak[4]
\[OM=\frac{OK}{\sin60^\circ}
=\frac{\frac32}{\frac{\sqrt3}{2}}=\sqrt3.\]
\Needspace{5\baselineskip}
于是
\nopagebreak[4]
\[\frac{\sqrt3R}{2}=\sqrt3\Longrightarrow R=2.\]
\Needspace{5\baselineskip}
球 $O$ 的表面积为
\nopagebreak[4]
\[S=4\pi R^2=4\pi\times4=16\pi.\]

\gkanswer{答案：$16\pi$。}

\end{answer}
\begin{answer}
\gksolhead{17}{}

\textbf{题目}：

等差数列 $\{a_n\}$ 的前 $n$ 项和为 $S_n$，已知 $S_3=a_2^2$，且 $S_1$，$S_2$，$S_4$ 成等比数列，求 $\{a_n\}$ 的通项公式。

\Needspace{4\baselineskip}
\textbf{解析}
\Needspace{5\baselineskip}
设等差数列的首项为 $a_1=a$，公差为 $d$，则
\nopagebreak[4]
\[S_3=3a+3d=3(a+d)=3a_2.\]
\Needspace{5\baselineskip}
由 $S_3=a_2^2$ 得
\nopagebreak[4]
\[3a_2=a_2^2,\]
\Needspace{5\baselineskip}
故
\nopagebreak[4]
\[a_2=0\quad\text{\text{或}}\quad a_2=3.\]

\Needspace{5\baselineskip}
\textbf{情形一：$a_2=3$。} 此时 $a+d=3$，即 $d=3-a$。于是
\nopagebreak[4]
\[S_1=a,\qquad S_2=2a+d=a+3,\]
\nopagebreak[4]
\[S_4=4a+6d=4a+6(3-a)=18-2a.\]
\Needspace{5\baselineskip}
由 $S_1$，$S_2$，$S_4$ 成等比数列得
\nopagebreak[4]
\[S_2^2=S_1S_4,\]
\Needspace{5\baselineskip}
即
\nopagebreak[4]
\[(a+3)^2=a(18-2a).\]
\Needspace{5\baselineskip}
展开：
\nopagebreak[4]
\[a^2+6a+9=18a-2a^2
\Longrightarrow 3a^2-12a+9=0
\Longrightarrow a^2-4a+3=0,\]
\Needspace{5\baselineskip}
所以
\nopagebreak[4]
\[a=1\quad\text{\text{或}}\quad a=3.\]
\Needspace{5\baselineskip}
若 $a=1$，则 $d=2$，
\nopagebreak[4]
\[a_n=1+2(n-1)=2n-1.\]
\Needspace{5\baselineskip}
若 $a=3$，则 $d=0$，
\nopagebreak[4]
\[a_n=3.\]
两者都使 $S_1$，$S_2$，$S_4$ 成等比数列（前者公比为 $4$，后者公比为 $2$）。

\Needspace{5\baselineskip}
\textbf{情形二：$a_2=0$。} 此时 $d=-a$，于是
\nopagebreak[4]
\[S_1=a,\qquad S_2=2a+d=a,\qquad S_4=4a+6d=-2a.\]
\Needspace{5\baselineskip}
由 $S_2^2=S_1S_4$ 得
\nopagebreak[4]
\[a^2=-2a^2\Longrightarrow 3a^2=0\Longrightarrow a=0.\]
此时 $S_1=S_2=S_4=0$，零数列不能作为等比数列（公比无定义），故舍去。

\Needspace{5\baselineskip}
综上，
\nopagebreak[4]
\[a_n=2n-1\quad\text{\text{或}}\quad a_n=3\quad(n\in\mathbf N^*).\]

\gkanswer{答案：$a_n=2n-1$ 或 $a_n=3$（$n\in\mathbf N^*$）。}

\end{answer}
\begin{answer}
\gksolhead{18}{}

\textbf{题目}：

设 $\triangle ABC$ 的内角 $A$，$B$，$C$ 的对边分别为 $a$，$b$，$c$，$(a+b+c)(a-b+c)=ac$。

（1） 求 $B$；

（2） 若 $\sin A\sin C=\dfrac{\sqrt3-1}4$，求 $C$。

\Needspace{4\baselineskip}
\textbf{解析}
\Needspace{5\baselineskip}
\textbf{（1）} 展开条件左边：
\nopagebreak[4]
\[\begin{aligned}
(a+b+c)(a-b+c)
&=(a+c)^2-b^2\\
&=a^2+2ac+c^2-b^2.
\end{aligned}\]
\Needspace{5\baselineskip}
由条件它等于 $ac$，故
\nopagebreak[4]
\[a^2+c^2-b^2=-ac.\]
\Needspace{5\baselineskip}
于是由余弦定理
\nopagebreak[4]
\[\cos B=\frac{a^2+c^2-b^2}{2ac}
=\frac{-ac}{2ac}=-\frac12.\]
\Needspace{5\baselineskip}
因 $0<B<\pi$，故
\nopagebreak[4]
\[B=\frac{2\pi}3.\]

\Needspace{5\baselineskip}
\textbf{（2）} 由（1），
\nopagebreak[4]
\[A+C=\pi-B=\frac\pi3.\]
\Needspace{5\baselineskip}
利用积化和差，
\nopagebreak[4]
\[\sin A\sin C
=\frac{\cos(A-C)-\cos(A+C)}2
=\frac{\cos(A-C)-\frac12}{2}.\]
\Needspace{5\baselineskip}
由条件
\nopagebreak[4]
\[\frac{\cos(A-C)-\frac12}{2}=\frac{\sqrt3-1}{4},\]
\Needspace{5\baselineskip}
故
\nopagebreak[4]
\[\cos(A-C)-\frac12=\frac{\sqrt3-1}{2}
\Longrightarrow \cos(A-C)=\frac{\sqrt3}{2}.\]
\Needspace{5\baselineskip}
因为 $A>0$，$C>0$ 且 $A+C=\dfrac\pi3$，所以
\nopagebreak[4]
\[-\frac\pi3<A-C<\frac\pi3.\]
\Needspace{5\baselineskip}
在此范围内 $\cos(A-C)=\dfrac{\sqrt3}{2}$ 给出
\nopagebreak[4]
\[A-C=\pm\frac\pi6.\]
\Needspace{5\baselineskip}
若 $A-C=\dfrac\pi6$，联立 $A+C=\dfrac\pi3$ 得
\nopagebreak[4]
\[A=\frac\pi4,\qquad C=\frac{\pi}{12}.\]
\Needspace{5\baselineskip}
若 $A-C=-\dfrac\pi6$，则
\nopagebreak[4]
\[C=\frac\pi4.\]
\Needspace{5\baselineskip}
故
\nopagebreak[4]
\[C=\frac{\pi}{12}\quad\text{\text{或}}\quad C=\frac\pi4.\]

\gkanswer{答案（1）：$B=\dfrac{2\pi}3$。}

\gkanswer{答案（2）：$C=\dfrac{\pi}{12}$ 或 $C=\dfrac\pi4$。}

\end{answer}
\begin{answer}
\gksolhead{19}{}

\textbf{题目}：

如图，四棱锥 $P\text{-}ABCD$ 中，$\angle ABC=\angle BAD=90^\circ$，$BC=2AD$，$\triangle PAB$ 与 $\triangle PAD$ 都是等边三角形。

（1） 证明：$PB\perp CD$；

（2） 求二面角 $A\text{-}PD\text{-}C$ 的大小。
\nopagebreak[4]
\begin{center}% 2013 年全国大纲卷理科第 19 题：四棱锥 P-ABCD。
% 来源 PDF：普通高考/2013/2013大纲理(广西).pdf，第 1 页（印面 795）右栏。
% 底面为直角梯形，AD∥BC，∠ABC=∠BAD=90°；CD、DA、PD 被遮挡。
\begin{tikzpicture}[semithick, font=\small, line join=round]
  \coordinate (C) at (0.8,1.0);
  \coordinate (B) at (8.2,1.0);
  \coordinate (A) at (9.4,2.6);
  \coordinate (D) at (5.6,2.6);
  \coordinate (P) at (6.2,6.4);
  % 可见棱
  \draw (P) -- (C);
  \draw (P) -- (B);
  \draw (P) -- (A);
  \draw (C) -- (B) -- (A);
  % 被遮挡的棱
  \draw[dashed] (P) -- (D);
  \draw[dashed] (C) -- (D) -- (A);
  % 标注
  \node[above] at (P) {$P$};
  \node[right] at (A) {$A$};
  \node[below] at (B) {$B$};
  \node[below left] at (C) {$C$};
  \node[above left] at (D) {$D$};
\end{tikzpicture}
\end{center}
\nopagebreak[4]
\Needspace{4\baselineskip}
\textbf{解析}
\Needspace{5\baselineskip}
\textbf{（1）} 因为 $\angle ABC=\angle BAD=90^\circ$，所以 $AB\perp BC$，$AB\perp AD$，从而 $AD\parallel BC$。由 $\triangle PAB$ 与 $\triangle PAD$ 都是等边三角形且共边 $PA$，得
\nopagebreak[4]
\[AB=AD=PA.\]
\Needspace{5\baselineskip}
又 $BC=2AD$。几何结论与整体尺度无关，不妨设
\nopagebreak[4]
\[AD=AB=2,\qquad BC=4.\]

\Needspace{5\baselineskip}
以 $A$ 为原点，$AB$，$AD$ 所在直线分别为 $x$，$y$ 轴，过 $A$ 垂直于底面的直线为 $z$ 轴建立坐标系：
\nopagebreak[4]
\[A(0,0,0),\quad B(2,0,0),\quad D(0,2,0).\]
\Needspace{5\baselineskip}
由 $AD\parallel BC$ 且 $BC=4$，结合图形中 $C$，$D$ 在 $AB$ 同侧，得
\nopagebreak[4]
\[C(2,4,0).\]
\Needspace{5\baselineskip}
设 $P(x,y,z)$。由 $PA=PB=PD=2$：
\nopagebreak[4]
\[x^2+y^2+z^2=4,\]
\nopagebreak[4]
\[(x-2)^2+y^2+z^2=4\Longrightarrow x=1,\]
\nopagebreak[4]
\[x^2+(y-2)^2+z^2=4\Longrightarrow y=1.\]
\Needspace{5\baselineskip}
于是 $z^2=4-1-1=2$，取 $z=\sqrt2$，得
\nopagebreak[4]
\[P(1,1,\sqrt2).\]
（取 $z=-\sqrt2$ 只把图形关于底面对称，不影响以下结论。）

\Needspace{5\baselineskip}
计算
\nopagebreak[4]
\[\overrightarrow{PB}=(1,-1,-\sqrt2),\qquad
\overrightarrow{CD}=(-2,-2,0).\]
\Needspace{5\baselineskip}
则
\nopagebreak[4]
\[\overrightarrow{PB}\cdot\overrightarrow{CD}
=1\times(-2)+(-1)\times(-2)+0=0.\]
故 $PB\perp CD$。

\Needspace{5\baselineskip}
\textbf{（2）} 平面 $APD$ 内两向量
\nopagebreak[4]
\[\overrightarrow{AP}=(1,1,\sqrt2),\qquad
\overrightarrow{AD}=(0,2,0).\]
\Needspace{5\baselineskip}
其法向量可取
\nopagebreak[4]
\[\boldsymbol n_1=\overrightarrow{AP}\times\overrightarrow{AD}
=(-2\sqrt2,0,2),\]
\Needspace{5\baselineskip}
化简为
\nopagebreak[4]
\[\boldsymbol n_1=(-\sqrt2,0,1).\]
\Needspace{5\baselineskip}
平面 $CPD$ 内两向量
\nopagebreak[4]
\[\overrightarrow{DC}=(2,2,0),\qquad
\overrightarrow{DP}=(1,-1,\sqrt2).\]
\Needspace{5\baselineskip}
其法向量可取
\nopagebreak[4]
\[\boldsymbol n_2=\overrightarrow{DC}\times\overrightarrow{DP}
=(2\sqrt2,-2\sqrt2,-4),\]
\Needspace{5\baselineskip}
化简为
\nopagebreak[4]
\[\boldsymbol n_2=(\sqrt2,-\sqrt2,-2).\]
\Needspace{5\baselineskip}
两法向量的夹角余弦的绝对值为
\nopagebreak[4]
\[\frac{|\boldsymbol n_1\cdot\boldsymbol n_2|}{|\boldsymbol n_1||\boldsymbol n_2|}
=\frac{4}{\sqrt3\cdot2\sqrt2}
=\frac{\sqrt6}{3}.\]
\Needspace{5\baselineskip}
二面角与两法向量的夹角相等或互补，下面用投影确定符号。取棱 $PD$ 的方向单位向量
\nopagebreak[4]
\[\boldsymbol u=\frac{\overrightarrow{PD}}{|\overrightarrow{PD}|}
=\left(-\frac12,\frac12,-\frac{\sqrt2}{2}\right).\]
\Needspace{5\baselineskip}
把 $\overrightarrow{PA}$，$\overrightarrow{PC}$ 投影到垂直于 $PD$ 的方向：
\nopagebreak[4]
\[\overrightarrow{PA}=(-1,-1,-\sqrt2),\qquad
\overrightarrow{PC}=(1,3,-\sqrt2).\]
\Needspace{5\baselineskip}
因为
\nopagebreak[4]
\[\overrightarrow{PA}\cdot\boldsymbol u=1,\qquad
\overrightarrow{PC}\cdot\boldsymbol u=2,\]
\Needspace{5\baselineskip}
所以
\nopagebreak[4]
\[\begin{aligned}
\boldsymbol v_A&=\overrightarrow{PA}-(\overrightarrow{PA}\cdot\boldsymbol u)\boldsymbol u
=\left(-\frac12,-\frac32,-\frac{\sqrt2}{2}\right),\\
\boldsymbol v_C&=\overrightarrow{PC}-(\overrightarrow{PC}\cdot\boldsymbol u)\boldsymbol u
=(2,2,0).
\end{aligned}\]
\Needspace{5\baselineskip}
这两个向量分别在两个半平面内且都垂直于棱 $PD$，把它们的起点平移到 $PD$ 上同一点，夹角不变，所得角即为二面角的平面角 $\theta$。于是
\nopagebreak[4]
\[\begin{aligned}
\cos\theta
&=\frac{\boldsymbol v_A\cdot\boldsymbol v_C}{|\boldsymbol v_A||\boldsymbol v_C|}\\
&=\frac{-1-3}{\sqrt3\cdot2\sqrt2}
=-\frac{4}{2\sqrt6}
=-\frac{\sqrt6}{3}.
\end{aligned}\]
\Needspace{5\baselineskip}
因此二面角为钝角，
\nopagebreak[4]
\[\theta=\arccos\left(-\frac{\sqrt6}{3}\right)
=\pi-\arccos\frac{\sqrt6}{3}
\approx144.7^\circ.\]

\gkanswer{答案（1）：证明见解析。}

\gkanswer{答案（2）：$\theta=\arccos\left(-\dfrac{\sqrt6}3\right)=\pi-\arccos\dfrac{\sqrt6}3$，约为 $144.7^\circ$。}

\end{answer}
\begin{answer}
\gksolhead{20}{}

\textbf{题目}：

甲、乙、丙三人进行羽毛球练习赛，其中两人比赛，另一人当裁判，每局比赛结束时，负的一方在下一局当裁判。设各局中双方获胜的概率均为 $\dfrac12$，各局比赛的结果都相互独立，第 $1$ 局甲当裁判。

（1） 求第 $4$ 局甲当裁判的概率；

（2） $X$ 表示前 $4$ 局中乙当裁判的次数，求 $X$ 的数学期望。

\Needspace{4\baselineskip}
\textbf{解析}
\Needspace{5\baselineskip}
设第 $n$ 局的裁判为 $J_n$。已知 $J_1=$ 甲。每局由另外两人比赛，负者担任下一局裁判，故 $J_{n+1}$ 是第 $n$ 局两名比赛者之一。由于第 $n$ 局的两名比赛者恰为除 $J_n$ 外的另外两人，且两人获胜概率各为 $\dfrac12$，所以
\nopagebreak[4]
\[P(J_{n+1}=X\mid J_n)=
\begin{cases}
\dfrac12,&X\neq J_n,\\
0,&X=J_n.
\end{cases}\]
也就是说，下一局裁判在另外两人中等可能选取。

\Needspace{5\baselineskip}
\textbf{（1）} 先求第 $3$ 局裁判的分布。由 $J_1=$ 甲，
\nopagebreak[4]
\[P(J_2=\text{\text{乙}})=P(J_2=\text{\text{丙}})=\frac12.\]
\Needspace{5\baselineskip}
若 $J_2=$ 乙，则 $J_3$ 等可能为甲或丙；若 $J_2=$ 丙，则 $J_3$ 等可能为甲或乙。故
\nopagebreak[4]
\[P(J_3=\text{\text{甲}})=\frac12,\qquad
P(J_3=\text{\text{乙}})=P(J_3=\text{\text{丙}})=\frac14.\]
\Needspace{5\baselineskip}
第 $4$ 局甲当裁判，必须 $J_3$ 为乙或丙，且由 $J_3$ 转移到甲的概率为 $\dfrac12$：
\nopagebreak[4]
\[\begin{aligned}
P(J_4=\text{\text{甲}})
&=P(J_3=\text{\text{乙}})\cdot\frac12
+P(J_3=\text{\text{丙}})\cdot\frac12\\
&=\frac14\cdot\frac12+\frac14\cdot\frac12=\frac14.
\end{aligned}\]

\Needspace{5\baselineskip}
\textbf{（2）} 第 $1$ 局甲当裁判，故 $X$ 是 $J_2,J_3,J_4$ 中“乙”出现的次数。每条长度为 $3$ 的裁判序列概率均为 $\left(\dfrac12\right)^3=\dfrac18$，共有 $8$ 条等可能的序列：
\nopagebreak[4]
\[\begin{array}{c|c}
(J_2,J_3,J_4)&X\\ \hline
(\text{\text{乙}},\text{\text{甲}},\text{\text{乙}})&2\\
(\text{\text{乙}},\text{\text{甲}},\text{\text{丙}})&1\\
(\text{\text{乙}},\text{\text{丙}},\text{\text{甲}})&1\\
(\text{\text{乙}},\text{\text{丙}},\text{\text{乙}})&2\\
(\text{\text{丙}},\text{\text{甲}},\text{\text{乙}})&1\\
(\text{\text{丙}},\text{\text{甲}},\text{\text{丙}})&0\\
(\text{\text{丙}},\text{\text{乙}},\text{\text{甲}})&1\\
(\text{\text{丙}},\text{\text{乙}},\text{\text{丙}})&1
\end{array}\]
\Needspace{5\baselineskip}
统计得
\nopagebreak[4]
\[P(X=0)=\frac18,\qquad
P(X=1)=\frac58,\qquad
P(X=2)=\frac28=\frac14.\]
\Needspace{5\baselineskip}
所以
\nopagebreak[4]
\[\begin{aligned}
E(X)&=0\cdot\frac18+1\cdot\frac58+2\cdot\frac14\\
&=\frac58+\frac12=\frac98.
\end{aligned}\]

\gkanswer{答案（1）：$\dfrac14$。}

\gkanswer{答案（2）：$P(X=0)=\dfrac18$，$P(X=1)=\dfrac58$，$P(X=2)=\dfrac14$，$E(X)=\dfrac98$。}

\end{answer}
\begin{answer}
\gksolhead{21}{}

\textbf{题目}：

已知双曲线 $C:\dfrac{x^2}{a^2}-\dfrac{y^2}{b^2}=1$（$a>0$，$b>0$）的左、右焦点分别为 $F_1$、$F_2$，离心率为 $3$，直线 $y=2$ 与 $C$ 的两个交点间的距离为 $\sqrt6$。

（1） 求 $a$、$b$；

（2） 设过 $F_2$ 的直线 $l$ 与 $C$ 的左、右两支分别交于 $A$、$B$ 两点，且 $|AF_1|=|BF_1|$，证明：$|AF_2|$、$|AB|$、$|BF_2|$ 成等比数列。

\Needspace{4\baselineskip}
\textbf{解析}
\Needspace{5\baselineskip}
\textbf{（1）} 离心率 $e=\dfrac ca=3$，故 $c=3a$，从而
\nopagebreak[4]
\[b^2=c^2-a^2=9a^2-a^2=8a^2.\]
\Needspace{5\baselineskip}
在 $y=2$ 处，
\nopagebreak[4]
\[\frac{x^2}{a^2}-\frac4{b^2}=1
\Longrightarrow x^2=a^2\left(1+\frac4{b^2}\right).\]
\Needspace{5\baselineskip}
两个交点关于 $y$ 轴对称，其距离为 $2|x|=\sqrt6$，故 $x^2=\dfrac32$。于是
\nopagebreak[4]
\[a^2\left(1+\frac4{b^2}\right)=\frac32.\]
\Needspace{5\baselineskip}
以 $b^2=8a^2$ 代入：
\nopagebreak[4]
\[a^2\left(1+\frac4{8a^2}\right)
=a^2+\frac12=\frac32,\]
\Needspace{5\baselineskip}
解得 $a^2=1$，即 $a=1$；进而
\nopagebreak[4]
\[b^2=8,\qquad b=2\sqrt2.\]

\Needspace{5\baselineskip}
\textbf{（2）} 由（1）得 $a=1$，$b=2\sqrt2$，$c=3$，故
\nopagebreak[4]
\[F_1(-3,0),\qquad F_2(3,0),\qquad |F_1F_2|=6.\]
\Needspace{5\baselineskip}
设
\nopagebreak[4]
\[|AF_1|=|BF_1|=t.\]
\Needspace{5\baselineskip}
$A$ 在左支上，由双曲线定义
\nopagebreak[4]
\[|AF_2|-|AF_1|=2a=2,\]
\Needspace{5\baselineskip}
故
\nopagebreak[4]
\[|AF_2|=t+2.\]
\Needspace{5\baselineskip}
$B$ 在右支上，由双曲线定义
\nopagebreak[4]
\[|BF_1|-|BF_2|=2a=2,\]
\Needspace{5\baselineskip}
故
\nopagebreak[4]
\[|BF_2|=t-2.\]
由 $|BF_2|>0$ 知 $t>2$。

\Needspace{5\baselineskip}
先排除 $F_2$ 在 $A$、$B$ 之间的情形。若 $F_2$ 在 $A$、$B$ 之间，设
\nopagebreak[4]
\[\angle AF_2F_1=\theta,\]
\Needspace{5\baselineskip}
则 $\angle BF_2F_1=\pi-\theta$。在 $\triangle AF_1F_2$ 与 $\triangle BF_1F_2$ 中分别用余弦定理：
\nopagebreak[4]
\[t^2=(t+2)^2+6^2-2\cdot6(t+2)\cos\theta,\]
\nopagebreak[4]
\[t^2=(t-2)^2+6^2-2\cdot6(t-2)\cos(\pi-\theta).\]
\Needspace{5\baselineskip}
第二式即
\nopagebreak[4]
\[t^2=(t-2)^2+36+12(t-2)\cos\theta.\]
\Needspace{5\baselineskip}
由第一式
\nopagebreak[4]
\[\cos\theta=\frac{(t+2)^2+36-t^2}{12(t+2)}
=\frac{t+10}{3(t+2)},\]
\Needspace{5\baselineskip}
由第二式
\nopagebreak[4]
\[\cos\theta=\frac{t^2-(t-2)^2-36}{12(t-2)}
=\frac{t-10}{3(t-2)}.\]
\Needspace{5\baselineskip}
两式相等：
\nopagebreak[4]
\[(t+10)(t-2)=(t-10)(t+2),\]
\Needspace{5\baselineskip}
展开得
\nopagebreak[4]
\[t^2+8t-20=t^2-8t-20
\Longrightarrow 16t=0,\]
与 $t>2$ 矛盾。故 $A$、$B$ 在 $F_2$ 的同侧。

\Needspace{5\baselineskip}
于是 $\angle AF_2F_1=\angle BF_2F_1=\theta$，且
\nopagebreak[4]
\[t^2=(t+2)^2+36-12(t+2)\cos\theta,\]
\nopagebreak[4]
\[t^2=(t-2)^2+36-12(t-2)\cos\theta.\]
\Needspace{5\baselineskip}
两式相减：
\nopagebreak[4]
\[0=8t-48\cos\theta
\Longrightarrow \cos\theta=\frac t6.\]
\Needspace{5\baselineskip}
代入第二式：
\nopagebreak[4]
\[\begin{aligned}
t^2&=(t-2)^2+36-12(t-2)\cdot\frac t6\\
&=t^2-4t+4+36-2t(t-2)\\
&=t^2-4t+40-2t^2+4t\\
&=40-t^2.
\end{aligned}\]
\Needspace{5\baselineskip}
故
\nopagebreak[4]
\[2t^2=40\Longrightarrow t=2\sqrt5.\]
\Needspace{5\baselineskip}
由于 $|AF_2|=t+2>t-2=|BF_2|$，$A$ 离 $F_2$ 更远，故 $F_2,B,A$ 依次共线，且
\nopagebreak[4]
\[|AB|=|AF_2|-|BF_2|=(t+2)-(t-2)=4.\]
\Needspace{5\baselineskip}
于是
\nopagebreak[4]
\[|AF_2||BF_2|=(t+2)(t-2)=t^2-4=20-4=16,\]
\Needspace{5\baselineskip}
而
\nopagebreak[4]
\[|AB|^2=4^2=16.\]
\Needspace{5\baselineskip}
所以
\nopagebreak[4]
\[|AB|^2=|AF_2||BF_2|,\]
即 $|AF_2|$、$|AB|$、$|BF_2|$ 成等比数列。

\gkanswer{答案（1）：$a=1$，$b=2\sqrt2$。}

\gkanswer{答案（2）：证明见解析。}

\end{answer}
\begin{answer}
\gksolhead{22}{}

\textbf{题目}：

已知函数 $f(x)=\ln(1+x)-\dfrac{x(1+\lambda x)}{1+x}$。

（1） 若 $x\geqslant0$ 时 $f(x)\leqslant0$，求 $\lambda$ 的最小值；

（2） 设数列 $\{a_n\}$ 的通项 $a_n=1+\dfrac12+\dfrac13+\cdots+\dfrac1n$，证明：$a_{2n}-a_n+\dfrac1{4n}>\ln2$。

\Needspace{4\baselineskip}
\textbf{解析}
\Needspace{5\baselineskip}
\textbf{（1）} 求导：
\nopagebreak[4]
\[\begin{aligned}
f'(x)&=\frac1{1+x}
-\frac{(1+2\lambda x)(1+x)-(x+\lambda x^2)}{(1+x)^2}\\
&=\frac{(1+x)-(1+2\lambda x+\lambda x^2)}{(1+x)^2}\\
&=\frac{x(1-2\lambda-\lambda x)}{(1+x)^2}.
\end{aligned}\]
因为 $f(0)=0$，要使 $x\geqslant0$ 时 $f(x)\leqslant0$，在 $x=0$ 右侧必须有 $f'(x)\leqslant0$。

\Needspace{5\baselineskip}
若 $\lambda<\dfrac12$，则 $1-2\lambda>0$，当 $x>0$ 充分小时
\nopagebreak[4]
\[1-2\lambda-\lambda x>0,\]
从而 $f'(x)>0$，于是 $f(x)>f(0)=0$，与条件矛盾。

\Needspace{5\baselineskip}
若 $\lambda\geqslant\dfrac12$，则 $1-2\lambda\leqslant0$，且 $-\lambda x\leqslant0$，故对一切 $x>0$ 有
\nopagebreak[4]
\[1-2\lambda-\lambda x\leqslant0,\]
\Needspace{5\baselineskip}
从而 $f'(x)\leqslant0$。因此 $f(x)$ 在 $[0,+\infty)$ 上单调不增，结合 $f(0)=0$ 得
\nopagebreak[4]
\[f(x)\leqslant0\quad(x\geqslant0).\]
\Needspace{5\baselineskip}
所以 $\lambda$ 的最小值为
\nopagebreak[4]
\[\lambda_{\min}=\frac12.\]
（当 $\lambda=\dfrac12$ 时，$f'(x)=-\dfrac{x^2}{2(1+x)^2}<0$，$x>0$。）

\Needspace{5\baselineskip}
\textbf{（2）} 由（1），取 $\lambda=\dfrac12$。当 $x>0$ 时 $f'(x)<0$，故 $f(x)<f(0)=0$，即
\nopagebreak[4]
\[\ln(1+x)<\frac{x\left(1+\frac x2\right)}{1+x}.\]
\Needspace{5\baselineskip}
对 $k=n+1,n+2,\dots,2n$，取 $x=\dfrac1{k-1}>0$，得
\nopagebreak[4]
\[\ln\frac{k}{k-1}
<\frac{\frac1{k-1}\left(1+\frac1{2(k-1)}\right)}{1+\frac1{k-1}}
=\frac{2k-1}{2k(k-1)}.\]
\Needspace{5\baselineskip}
而
\nopagebreak[4]
\[\frac{2k-1}{2k(k-1)}
=\frac1k+\frac1{2k(k-1)}
=\frac1k+\frac12\left(\frac1{k-1}-\frac1k\right).\]
\Needspace{5\baselineskip}
对 $k=n+1$ 到 $2n$ 求和：
\nopagebreak[4]
\[\begin{aligned}
\ln 2
&=\sum_{k=n+1}^{2n}\ln\frac{k}{k-1}\\
&<\sum_{k=n+1}^{2n}\frac1k
+\frac12\sum_{k=n+1}^{2n}\left(\frac1{k-1}-\frac1k\right)\\
&=(a_{2n}-a_n)
+\frac12\left(\frac1n-\frac1{2n}\right)\\
&=(a_{2n}-a_n)+\frac1{4n}.
\end{aligned}\]
\Needspace{5\baselineskip}
移项即得
\nopagebreak[4]
\[a_{2n}-a_n+\frac1{4n}>\ln2.\]

\gkanswer{答案（1）：$\lambda_{\min}=\dfrac12$。}

\gkanswer{答案（2）：证明见解析。}

\end{answer}